How do you setup your project to be most easily handled through F4PGA? 
Unlike with Vivado you'll need to handle managing the folder structure yourself.
It is a good idea to begin using git to start managing your code so that if you find yourself in a pickle you can revert back to functioning code.

So long as the Makefile directs to all the correct locations most structures can work; 
however, we're going to walk through the following structure:

\begin{figure}[H]
    \centering
    \begin{lstlisting}[language=tree]
    f4pga_project/
    ├── .git/
    │   ├── config
    │   ├── HEAD
    │   ├── objects/
    │   ├── refs/
    ├── src/
    │   ├── top.sv
    │   ├── module1.sv
    │   ├── module2.sv
    │   ├── module3.sv
    │   ├── ...
    ├── constraints/
    │   ├── constraints.xdc
    ├── build/
    │   ├── output.bit
    │   ├── log.txt
    ├── common/
    │   ├── common.mk
    ├── Makefile
    ├── README.md
    └── .gitignore
    \end{lstlisting}
    \caption{F4PGA Project Directory Structure}
\end{figure}



\begin{itemize}
    \item .git folder, You are going to want to manage your project using git. 
    This is a good idea regardless of toolchain; however, with f4pga it becomes especially important to ensure that directory strucutres don't get disturbed.
    If you're using github, you'll likely have a .github folder; however, this will for the most part be managed for you. 

    \item src folder, this will be where you keep all your modules not including your testbenches. 

    \item constraints folder, this is where you'll store your xdc file or multiple xdc files if your project is complicated enough to justify that.

    \item build folder, this should not be included in your git repository. 
    It is what your makefile will place its output into and where you'll find your output bit file.

    \item common folder, this can be included in your git repository.
    It contains the additional make instructions that will be combined with your makefile to complete the bitstream generation. 
    You may notice we didn't include this in the mount for docker. 
    This is technically true; however, 
    the container already contains a copy of the f4pga common makefile internally which ends up being referenced by the specified mount. 
    It is for the best to simply include it yourself to avoid any hassles regarding directory structure. 

    \item Makefile, we will go through the structure of this file in detail; however, it is what defines making the project. 
    If you paid attention earlier you may have noticed that we specified the basys 3 when building our project. 
    If you have multiple potential target boards, then you'll want multiple elements in your makefile showing how each board is compiled.

    \item README.md, document your project. This is important both for yourself and for others who may be looking at your code down the line. 
    This is a best practice worth getting used to!

    \item .gitignore. This tells git what not to include in the github repository. 
    In this case it'll be the build folder that you want to eliminate; however, in other cases you may eliminate things such as pyc files in python projects. 
    
\end{itemize}

\subsubsection{F4PGA Makefiles}
F4PGA has a good guide on creating custom makefiles for your projects, for a more extensive explanation go to: 
https://f4pga-examples.readthedocs.io/en/latest/customizing-makefiles.html.

\begin{lstlisting}[language=make]
    current_dir := ${CURDIR}

    TARGET := basys3

    TOP := <put the name of your top module here>

    SOURCES := ${current_dir}/src/*.sv

    XDC := ${current_dir}/constraints/constraints.xdc

    include ${current_dir}/common/common.mk
\end{lstlisting}


I've filled in this example with reasonable guesses for what each of your values should be in the makefile assuming you have identical folder structure to what is shown.
The idea is that this file specifies all the details of the design. 
common.mk is provided both in the repo for this document and in f4pga-examples repo. 
It contains all the required instructions to the toolchain to export a bitstream file. 